\documentclass[14pt,a4paper,oneside]{extbook}
\usepackage[utf8]{vietnam}
\usepackage{amsmath,amssymb,mathrsfs,fancyhdr,enumerate,multirow,makecell}
\usepackage{tkz-euclide,tikz-3dplot,tikz,tkz-tab}
\usetikzlibrary{arrows.meta,calc,intersections,angles,quotes,backgrounds,shapes.geometric,patterns,through,shadows,positioning}
\usepackage{yhmath}
\usepackage[top=0.7cm, bottom=0.7cm, left=1.0cm, right=1.0cm, headheight=16pt, headsep=3mm, footskip=5mm]{geometry}
\usepackage[hidelinks,unicode]{hyperref}
\usepackage{qrcode}
\usepackage[many]{tcolorbox}
\usepackage{fontawesome}

\pagestyle{fancy}
\fancyhf{}
\fancyhead[L]{\footnotesize\bfseries\color{teal!85!black}\faBook~SỔ TAY KIẾN THỨC TOÁN 4}
\fancyhead[R]{\footnotesize\bfseries\color{blue!85!black}Lớp Toán Thầy Lĩnh -- loptoanthaylinh.vn}
\fancyfoot[C]{\footnotesize\bfseries\color{gray!80!black}-- Sổ tay kiến thức: Bài 19. Giây, thế kỉ --}
\renewcommand{\headrulewidth}{0.4pt}

\begin{document}

% --- BANNER TIÊU ĐỀ BÀI HỌC ---
\noindent
\begin{tcolorbox}[
  enhanced,
  colback=teal!8!white,
  colframe=teal!80!black,
  arc=3mm, boxrule=1pt,
  left=3mm, right=3mm, top=2mm, bottom=2mm,
  halign=center
]
  {\Large\bfseries\color{teal!90!black}\faBookmark\ BÀI 19: GIÂY, THẾ KỈ}\\[1.5pt]
  {\footnotesize\itshape\color{black!75} Bảng đơn vị thời gian giây và thế kỉ, chữ số La Mã và phương pháp xác định thế kỉ trên 1 trang A4}
\end{tcolorbox}

\vspace{2mm}

% --- KHỐI 1 & 2: LÝ THUYẾT & KĨ NĂNG ---
\noindent
\begin{minipage}[t]{0.49\textwidth}
  \begin{tcolorbox}[
    enhanced,
    title={\footnotesize\bfseries\color{white}\faClockO\ 1. Đơn vị Giây \& Đổi thời gian},
    colback=blue!2!white,
    colframe=blue!75!black,
    arc=2.5mm, boxrule=0.8pt, top=1.8mm, bottom=1.8mm, left=2mm, right=2mm
  ]
    \footnotesize
    $\bullet$ \textbf{Giây} là đơn vị đo thời gian bé hơn phút:\\
    --- $1\text{ phút} = 60\text{ giây}$\\
    --- $1\text{ giờ} = 60\text{ phút} = 3\,600\text{ giây}$\\
    $\bullet$ \textbf{Một số phân số thời gian thường gặp}:\\
    --- $\dfrac{1}{2}\text{ phút} = 30\text{ giây}$\\
    --- $\dfrac{1}{4}\text{ phút} = 15\text{ giây}$\\
    --- $\dfrac{1}{3}\text{ phút} = 20\text{ giây}$
  \end{tcolorbox}
\end{minipage}%
\hfill
\begin{minipage}[t]{0.49\textwidth}
  \begin{tcolorbox}[
    enhanced,
    title={\footnotesize\bfseries\color{white}\faHistory\ 2. Thế kỉ \& Cách xác định thế kỉ},
    colback=teal!2!white,
    colframe=teal!75!black,
    arc=2.5mm, boxrule=0.8pt, top=1.8mm, bottom=1.8mm, left=2mm, right=2mm
  ]
    \footnotesize
    $\bullet$ \textbf{Thế kỉ} là đơn vị đo thời gian: $1\text{ thế kỉ} = 100\text{ năm}$. Viết bằng chữ số La Mã: I, II, ..., XX, XXI.\\
    $\bullet$ \textbf{Quy tắc xác định năm thuộc thế kỉ nào}:\\
    --- Nếu năm tận cùng là \textbf{00}: Lấy số phía trước hai chữ số 0.\\
    \textit{VD:} Năm $2000 \rightarrow$ Thế kỉ \textbf{XX} (thế kỉ 20).\\
    --- Nếu năm \textbf{không} tận cùng là 00: Lấy số phía trước hai chữ số cuối rồi \textbf{cộng thêm 1}.\\
    \textit{VD:} Năm $1945 \rightarrow 19 + 1 = 20 \Rightarrow$ Thế kỉ \textbf{XX}.
  \end{tcolorbox}
\end{minipage}

\vspace{2mm}

% --- KHỐI 3: SƠ ĐỒ TƯ DUY TIKZ ---
\noindent
\begin{tcolorbox}[
  enhanced,
  title={\footnotesize\bfseries\color{white}\faSitemap\ SƠ ĐỒ TƯ DUY: TRỤC LỊCH SỬ CÁC THẾ KỈ VÀ MẸO TÍNH NHANH THẾ KỈ},
  colback=yellow!2!white,
  colframe=orange!85!black,
  arc=2.5mm, boxrule=0.8pt, top=1.5mm, bottom=1.5mm, left=1mm, right=1mm,
  halign=center
]
  \begin{tikzpicture}[
    xscale=1.1, yscale=0.85,
    tag/.style={rectangle, rounded corners=1.5mm, draw=teal!80!black, fill=teal!10, inner sep=2.5pt, font=\scriptsize\bfseries, align=center},
    arr/.style={-{Stealth[scale=0.8]}, thick, draw=blue!80!black}
  ]
    % Trục thế kỉ
    \draw[-{Stealth[scale=1.2]}, very thick, draw=teal!85!black] (-0.3, 0) -- (12.2, 0);

    \node[tag, fill=blue!12, draw=blue!80!black] at (1.2, 0.6) {Thế kỉ I\\Năm 1 -- 100};
    \node[tag, fill=teal!12, draw=teal!80!black] at (4.0, 0.6) {Thế kỉ XIX\\Năm 1801 -- 1900};
    \node[tag, fill=orange!15, draw=orange!85!black] at (7.3, 0.6) {Thế kỉ XX\\Năm 1901 -- 2000};
    \node[tag, fill=purple!15, draw=purple!85!black] at (10.6, 0.6) {Thế kỉ XXI\\Năm 2001 -- 2100};

    \draw[thick, draw=black!70] (1.2, 0.2) -- (1.2, -0.2);
    \draw[thick, draw=black!70] (4.0, 0.2) -- (4.0, -0.2);
    \draw[thick, draw=black!70] (7.3, 0.2) -- (7.3, -0.2);
    \draw[thick, draw=black!70] (10.6, 0.2) -- (10.6, -0.2);

    % Khối chú thích phía dưới
    \node[rectangle, rounded corners=2mm, draw=red!70!black, fill=red!8, inner sep=3pt, font=\scriptsize, align=left] at (3.5, -0.9) {
      \textbf{Tận cùng là 00:} Bỏ hai số 0 $\rightarrow$ Chính là thế kỉ!\\
      \textit{Ví dụ:} $1500 \rightarrow$ Thế kỉ XV; $1900 \rightarrow$ Thế kỉ XIX.
    };
    \node[rectangle, rounded corners=2mm, draw=green!70!black, fill=green!8, inner sep=3pt, font=\scriptsize, align=left] at (9.2, -0.9) {
      \textbf{Tận cùng khác 00:} Bỏ hai số cuối rồi $\mathbf{+ 1}$\\
      \textit{Ví dụ:} $2024 \rightarrow 20 + 1 = 21 \Rightarrow$ Thế kỉ XXI.
    };
  \end{tikzpicture}
\end{tcolorbox}

\vspace{2mm}

% --- KHỐI 4: VÍ DỤ MINH HỌA & BẪY TRẮC NGHIỆM ---
\noindent
\begin{minipage}[t]{0.53\textwidth}
  \begin{tcolorbox}[
    enhanced,
    title={\footnotesize\bfseries\color{white}\faCheckCircle\ Ví dụ thực hành chuẩn mực},
    colback=green!2!white,
    colframe=green!60!black,
    arc=2.5mm, boxrule=0.8pt, top=1.8mm, bottom=1.8mm, left=2.2mm, right=2.2mm
  ]
    \footnotesize
    $\bullet$ \textbf{Ví dụ 1 (Đổi thời gian):}\\
    $2\text{ phút } 15\text{ giây} = 2 \times 60 + 15 = \mathbf{135\text{ giây}}$.\\
    $200\text{ giây} = 180\text{ giây} + 20\text{ giây} = \mathbf{3\text{ phút } 20\text{ giây}}$.\\
    $\bullet$ \textbf{Ví dụ 2 (Lịch sử):}\\
    --- Bác Hồ đọc Tuyên ngôn Độc lập năm $1945$: Thuộc thế kỉ \textbf{XX} (vì $19 + 1 = 20$).\\
    --- Chiến thắng Điện Biên Phủ năm $1954$: Thế kỉ \textbf{XX}.\\
    --- Kỉ niệm 1000 năm Thăng Long năm $2010$: Thế kỉ \textbf{XXI}.
  \end{tcolorbox}
\end{minipage}%
\hfill
\begin{minipage}[t]{0.45\textwidth}
  \begin{tcolorbox}[
    enhanced,
    title={\footnotesize\bfseries\color{white}\faExclamationTriangle\ Mẹo tránh bẫy trắc nghiệm},
    colback=red!2!white,
    colframe=red!70!black,
    arc=2.5mm, boxrule=0.8pt, top=1.8mm, bottom=1.8mm, left=2.2mm, right=2.2mm
  ]
    \footnotesize
    $\bullet$ \textbf{Bẫy năm chẵn trăm 2000:}\\
    Nhiều bạn thấy số 20 nên nghĩ năm 2000 là thế kỉ 21 $\rightarrow$ \textbf{SAI}! Năm 2000 là năm kết thúc của thế kỉ \textbf{XX}.\\
    $\bullet$ \textbf{Bẫy chữ số La Mã:}\\
    Nhớ chính xác các số La Mã: $X = 10, IX = 9, XIV = 14, XIX = 19, XX = 20, XXI = 21$.
  \end{tcolorbox}
\end{minipage}

\end{document}
